\section{} % 1

\subsection{} % 1.1

We write the translate-rotate-scale function \lstinline{trs()} as follows:

\begin{lstlisting}
def trs(
    image: NDArray[Any],
    t: tuple[float, float],
    theta: float,
    scale: float,
) -> NDArray[Any]:
    H, W = image.shape
    assert image.ndim == 2, "Input image must be 2D grayscale image"
    assert len(t) == 2, "t should be a tuple of length 2"
    assert scale > 0, "scale should be a positive number"
    assert H % 2 == 1, "Input image must have uneven height"
    assert W % 2 == 1, "Input image must have uneven width"
    
    c = (W // 2, H // 2)  # center of the image

    # We give the transformation product in its inverse form
    Tt_inv = _translate_t((-t[0], -t[1]))
    Tc_inv = _translate_t((-c[0], -c[1]))
	Tc = _translate_t(c)
    R_inv = _rotate(-theta)
    S_inv = np.diag([1/scale, 1/scale, 1])

    inv_transformation = Tc @ S_inv @ R_inv @ Tc_inv @ Tt_inv

    return warp(image, inv_transformation, order=0)
\end{lstlisting}
which makes use of the \lstinline{_translate_t()} and \lstinline{_rotate()} helper functions defined as follows:

\begin{lstlisting}
def _translate_t(t: tuple[float, ...], ndim: int = 2) -> NDArray[Any]:
    assert len(t) == ndim, "t should have the same number of dimensions as ndim"
    Tt = np.identity(ndim + 1)
    for i in range(len(t)):
        Tt[i, ndim] = t[i]

    return Tt
\end{lstlisting}
a translation function, that takes in a translation vector and returns the corresponding translation matrix
in homogeneous coordinates, and
\begin{lstlisting}
def _rotate(theta: float, ndim: int = 2) -> NDArray[Any]:
    """Returns the rotation matrix for a given angle theta in radians."""
    assert ndim == 2, "`_rotate` is only defined for 2D images"
    R = np.identity(ndim + 1)
    R[0, 0] = np.cos(theta)
    R[0, 1] = np.sin(theta)
    R[1, 0] = -np.sin(theta)
    R[1, 1] = np.cos(theta)

    return R
\end{lstlisting}
a rotation function that takes in an angle in radians and returns the corresponding rotation matrix in homogeneous coordinates.

One key implementation detail for the \lstinline{trs()} function is that we give the transformation product in its inverse form, because \lstinline{skimage.transform.warp} applies the inverse of the given transformation. Thus since $(AB)^{-1} = B^{-1}A^{-1}$, we give the product of the inverse transformations in reverse order i.e. $(T_tT_cRST_c^{-1})^{-1}=T_cT_c^{-1}R^{-1}S^{-1}T_t^{-1}$.
For translation the inverse translation is simply the negative of the translation vector.
For rotation, the inverse is the negative of the angle. 
For scaling, the inverse is the reciprocal of the scale factor.
The inverses thus only depend on the input arguments, which makes the implementation of the \lstinline{trs()} function straightforward.
It is important to set the keyword argument \lstinline{order=0} when calling \lstinline{skimage.transform.warp} to use nearest neighbor interpolation,
since the default interpolation method is bilinear interpolation.

Below in figure \ref{fig:trs_square} we show the result of applying the \lstinline{trs()} function to the white square image:

\begin{figure}[H]
	\centering
	\includegraphics[width=0.8\textwidth]{ass8/transformed_image.png}
	\caption{
		The result of applying the \lstinline{trs()} function to the white square image from Q.14 Assignment 6 Week 6.
		The transformation parameters are $s=2$, $\theta = \pi/10$, and $t=(10.4, 15.7)$.
	}
	\label{fig:trs_square}
\end{figure}

Judging from the figure \ref{fig:trs_square}, by visual inspection the transformation seems to be applied correctly.
Clearly the rotation is counter-clockwise about $\pi/10$ radians which is the same as $18^\circ$.
The scaling also seems about double. The translation occuring is not as simple to visually verify,
but the center of the square does indeed seem to be translated.
